% Общая характеристика работы --- двуязычная (переключатель thesislang).
% Используем примитивный \ifnum...\else...\fi, чтобы не ломать \verb внутри.
\ifnum\value{thesislang}=1
%%%%%%%%%%%%%%%%%%%%%%%%%%%%%%%%% English placeholder %%%%%%%%%%%%%%%%%%%%%%%%%%%
% The guidance in the comments below is adapted from O.\,V.~Vasilyev,
% "Preparing a Dissertation for the Degree of Candidate of Sciences:
% a 10-Minute Guide" (Skoltech). The comments (%) do not appear in the PDF ---
% they are here to help you write each section. Replace every \fixme{...}.

% A brief literature review focused on the *significance* of the research area:
% recent progress, impact on science and/or industry. Discuss what is already
% addressed (with references) and highlight what is not (or only partially)
% addressed, then articulate why now is the right time for this research and how
% it fits the field's agenda.
{\actualityandprogress} \fixme{Introduce the topic, its relevance and the place of
this work in the field. You may cite other works with \autocite{Panov2025}.}

% The goal must be *concrete* --- to develop a theory / model / methodology, to
% demonstrate the applicability of one, etc. "To study ..." is NOT a goal.
% Note: "цели и задачи" = "goals and problems addressed" (not "aims/objectives/
% tasks"). Problems are specific (technical, theoretical, experimental) and are
% distinct from the broader "main results submitted for the defense" below.
{\goalsproblems} The {\goal} of this dissertation is to \fixme{develop \dots}

To achieve the goal of the dissertation, the following {\tasks} are addressed:
\begin{enumerate}
  \item \fixme{To develop a \dots{} (model, method, methodology, \dots).}
  \item \fixme{To verify and validate the developed \dots{} against \dots}
  \item \fixme{To apply the developed \dots{} to \dots}
\end{enumerate}

% Focus on what is *novel*, not on the problems solved: a new model for an
% unstudied behaviour, new regimes described, a theory generalized, etc.
% Useful phrasings:
%   - The thesis presents a fundamentally new (method, approach, ...) with unique
%     properties that distinguish it from all known ...
%   - The dissertation also presents the novel (model, method, theory, ...)
%     developed by the author ...
%   - The author's groundbreaking results also include ...
{\novelty}
\begin{enumerate}
  \item \fixme{The thesis presents a fundamentally new \dots{} with unique
  properties that distinguish it from all known \dots}
  \item \fixme{The dissertation also presents the novel \dots{} developed by the
  author.}
  \item \fixme{The author's groundbreaking results also include \dots}
\end{enumerate}

% Which results have theoretical and/or practical significance, and why. If a
% new approach founds a research direction, say how it may develop; if results
% apply to other fields, list them; explain the significance in plain terms.
{\theoreticalsignificance} \fixme{of the results of the dissertation lies in the
development of \dots{} which lays the foundation for \dots}

{\practicalsignificance} \fixme{of this study lies in \dots{}, which enables
\dots}

% Focus on the *tools* --- methods, theories, approaches, experimental
% procedures or facilities the research is based on or used to verify/validate
% the results --- NOT on the author's achievements.
{\methods}
\fixme{The study is based on \dots{} (methods, theories, approaches). The results
are verified/validated using \dots{} (experimental procedures or facilities).}

% Of primary importance: these are read verbatim at the defense and the defense
% presentation is built around them. Distinct from "goals and problems": the
% main results are grandiose and assessable (a model/theory developed, a theorem
% proved, a method proposed, a phenomenon discovered) and must be unique to you.
{\defpositions}
\begin{enumerate}
  \item \fixme{A \dots{} is developed (a model/theory developed, a theorem
  proved, a method proposed, a new phenomenon discovered, \dots).}
  \item \fixme{\dots}
\end{enumerate}

% Clarify your contribution to each paper/chapter the dissertation is based on,
% especially for multi-author papers; substantiate its nature and, where
% applicable, emphasize that it is defining. Write in the third person.
{\contribution} \fixme{The author had a decisive contribution to \dots}

% Validation is broader than "simulation vs. experiment" --- include
% reproducibility, detailed method descriptions, agreement with published
% results, etc. Typical opening: "The validity of the results of the
% dissertation is confirmed by the following."
{\reliability} \fixme{The validity of the results of the dissertation is confirmed
by the following: the developed methods are described in detail in \dots; they are
implemented in \dots; the results are in agreement with published experimental
and/or computational results obtained by \dots}

% List the conferences/venues where the results were presented (and, together
% with {\publications} below, the publications). May be merged with validity.
{\probation} The main results of the work were presented~at:
\begin{itemize}
  \item \fixme{Conference name, year, city, country.}
\end{itemize}

% Suggested phrasing: "The main results on the topic of the dissertation are
% presented in N publications, M of which in journals indexed by Web of Science
% or Scopus, K in conference abstracts."
{\publications} \fixme{The main results on the topic of the dissertation are
presented in N~publications, M~of which in journals indexed by Web of Science or
Scopus, and K~in conference abstracts.}
%%%%%%%%%%%%%%%%%%%%%%%%%%%%%%%%%%%%%%%%%%%%%%%%%%%%%%%%%%%%%%%%%%%%%%%%%%%%%%%%%%
\else
%%%%%%%%%%%%%%%%%%%%%%%%%%%%%%%%% Русский вариант %%%%%%%%%%%%%%%%%%%%%%%%%%%%%%%%
{\actuality} Обзор, введение в тему, обозначение места данной работы в
мировых исследованиях и~т.\:п., можно использовать ссылки на~другие
работы~\autocite{Panov2025}
(если их~нет, то~в~автореферате
автоматически пропадёт раздел <<Список литературы>>). Внимание! Ссылки
на~другие работы в~разделе общей характеристики работы можно
использовать только при использовании \verb!biblatex! (из-за технических
ограничений \verb!bibtex8!. Это связано с тем, что одна
и~та~же~характеристика используются и~в~тексте диссертации, и в
автореферате. В~последнем, согласно ГОСТ, должен присутствовать список
работ автора по~теме диссертации, а~\verb!bibtex8! не~умеет выводить в~одном
файле два списка литературы).
При использовании \verb!biblatex! возможно использование исключительно
в~автореферате подстрочных ссылок
для других работ командой \verb!\autocite!~\autocite{Panov2025}, а~также цитирование
собственных работ командой \verb!\cite!. Для этого в~файле
\verb!common/setup.tex! необходимо присвоить положительное значение
счётчику \verb!\setcounter{usefootcite}{1}!.

Для генерации содержимого титульного листа автореферата, диссертации
и~презентации используются данные из файла \verb!common/data.tex!. Если,
например, вы меняете название диссертации, то оно автоматически
появится в~итоговых файлах после очередного запуска \LaTeX. Согласно
ГОСТ 7.0.11-2011 <<5.1.1 Титульный лист является первой страницей
диссертации, служит источником информации, необходимой для обработки и
поиска документа>>. Наличие логотипа организации на~титульном листе
упрощает обработку и~поиск, для этого разметите логотип вашей
организации в папке images в~формате PDF (лучше найти его в векторном
варианте, чтобы он хорошо смотрелся при печати) под именем
\verb!logo.pdf!. Настроить размер изображения с логотипом можно
в~соответствующих местах файлов \verb!title.tex!  отдельно для
диссертации и автореферата. Если вам логотип не~нужен, то просто
удалите файл с~логотипом.

\ifsynopsis
Этот абзац появляется только в~автореферате.
Для формирования блоков, которые будут обрабатываться только в~автореферате,
заведена проверка условия \verb!\!\verb!ifsynopsis!.
Значение условия задаётся в~основном файле документа (\verb!synopsis.tex! для
автореферата).
\else
Этот абзац появляется только в~диссертации.
Через проверку условия \verb!\!\verb!ifsynopsis!, задаваемого в~основном файле
документа (\verb!dissertation.tex! для диссертации), можно сделать новую
команду, обеспечивающую появление цитаты в~диссертации, но~не~в~автореферате.
\fi

% {\progress}
% Этот раздел должен быть отдельным структурным элементом по
% ГОСТ, но он, как правило, включается в описание актуальности
% темы. Нужен он отдельным структурынм элемементом или нет ---
% смотрите другие диссертации вашего совета, скорее всего не нужен.

{\aim} данной работы является \ldots

Для~достижения поставленной цели необходимо было решить следующие {\tasks}:
\begin{enumerate}[beginpenalty=10000] % https://tex.stackexchange.com/a/476052/104425
  \item Исследовать, разработать, вычислить и~т.\:д. и~т.\:п.
  \item Исследовать, разработать, вычислить и~т.\:д. и~т.\:п.
  \item Исследовать, разработать, вычислить и~т.\:д. и~т.\:п.
  \item Исследовать, разработать, вычислить и~т.\:д. и~т.\:п.
\end{enumerate}


{\novelty}
\begin{enumerate}[beginpenalty=10000] % https://tex.stackexchange.com/a/476052/104425
  \item Впервые \ldots
  \item Впервые \ldots
  \item Было выполнено оригинальное исследование \ldots
\end{enumerate}

{\influence} \ldots

{\methods} \ldots

{\defpositions}
\begin{enumerate}[beginpenalty=10000] % https://tex.stackexchange.com/a/476052/104425
  \item Первое положение
  \item Второе положение
  \item Третье положение
  \item Четвертое положение
\end{enumerate}
В папке Documents можно ознакомиться с решением совета из Томского~ГУ
(в~файле \verb+Def_positions.pdf+), где обоснованно даются рекомендации
по~формулировкам защищаемых положений.

{\reliability} полученных результатов обеспечивается \ldots \ Результаты находятся в соответствии с результатами, полученными другими авторами.


{\probation}
Основные результаты работы докладывались~на:
перечисление основных конференций, симпозиумов и~т.\:п.

{\contribution} Автор принимал активное участие \ldots

\ifnumequal{\value{bibliosel}}{0}
{%%% Встроенная реализация с загрузкой файла через движок bibtex8. (При желании, внутри можно использовать обычные ссылки, наподобие `\cite{vakbib1,vakbib2}`).
    {\publications} Основные результаты по теме диссертации изложены
    в~XX~печатных изданиях,
    X из которых изданы в журналах, рекомендованных ВАК,
    X "--- в тезисах докладов.
}%
{%%% Реализация пакетом biblatex через движок biber
    \begin{refsection}[bl-author, bl-registered]
        % Это refsection=1.
        % Процитированные здесь работы:
        %  * подсчитываются, для автоматического составления фразы "Основные результаты ..."
        %  * попадают в авторскую библиографию, при usefootcite==0 и стиле `\insertbiblioauthor` или `\insertbiblioauthorgrouped`
        %  * нумеруются там в зависимости от порядка команд `\printbibliography` в этом разделе.
        %  * при использовании `\insertbiblioauthorgrouped`, порядок команд `\printbibliography` в нём должен быть тем же (см. biblio/biblatex.tex)
        %
        % Невидимый библиографический список для подсчёта количества публикаций:
        \phantom{\printbibliography[heading=nobibheading, section=1, env=countauthorvak,          keyword=biblioauthorvak]%
        \printbibliography[heading=nobibheading, section=1, env=countauthorwos,          keyword=biblioauthorwos]%
        \printbibliography[heading=nobibheading, section=1, env=countauthorscopus,       keyword=biblioauthorscopus]%
        \printbibliography[heading=nobibheading, section=1, env=countauthorconf,         keyword=biblioauthorconf]%
        \printbibliography[heading=nobibheading, section=1, env=countauthorother,        keyword=biblioauthorother]%
        \printbibliography[heading=nobibheading, section=1, env=countregistered,         keyword=biblioregistered]%
        \printbibliography[heading=nobibheading, section=1, env=countauthorpatent,       keyword=biblioauthorpatent]%
        \printbibliography[heading=nobibheading, section=1, env=countauthorprogram,      keyword=biblioauthorprogram]%
        \printbibliography[heading=nobibheading, section=1, env=countauthor,             keyword=biblioauthor]%
        \printbibliography[heading=nobibheading, section=1, env=countauthorvakscopuswos, filter=vakscopuswos]%
        \printbibliography[heading=nobibheading, section=1, env=countauthorscopuswos,    filter=scopuswos]}%
        %
        \nocite{*}%
        %
        {\publications} Основные результаты по теме диссертации изложены в~\arabic{citeauthor}~печатных изданиях,
        \arabic{citeauthorvak} из которых изданы в журналах, рекомендованных ВАК%
        \ifnum \value{citeauthorscopuswos}>0%
            , \arabic{citeauthorscopuswos} "--- в~периодических научных журналах, индексируемых Web of~Science и Scopus%
        \fi%
        \ifnum \value{citeauthorconf}>0%
            , \arabic{citeauthorconf} "--- в~тезисах докладов.
        \else%
            .
        \fi%
        \ifnum \value{citeregistered}=1%
            \ifnum \value{citeauthorpatent}=1%
                Зарегистрирован \arabic{citeauthorpatent} патент.
            \fi%
            \ifnum \value{citeauthorprogram}=1%
                Зарегистрирована \arabic{citeauthorprogram} программа для ЭВМ.
            \fi%
        \fi%
        \ifnum \value{citeregistered}>1%
            Зарегистрированы\ %
            \ifnum \value{citeauthorpatent}>0%
            \formbytotal{citeauthorpatent}{патент}{}{а}{}%
            \ifnum \value{citeauthorprogram}=0 . \else \ и~\fi%
            \fi%
            \ifnum \value{citeauthorprogram}>0%
            \formbytotal{citeauthorprogram}{программ}{а}{ы}{} для ЭВМ.
            \fi%
        \fi%
        % К публикациям, в которых излагаются основные научные результаты диссертации на соискание учёной
        % степени, в рецензируемых изданиях приравниваются патенты на изобретения, патенты (свидетельства) на
        % полезную модель, патенты на промышленный образец, патенты на селекционные достижения, свидетельства
        % на программу для электронных вычислительных машин, базу данных, топологию интегральных микросхем,
        % зарегистрированные в установленном порядке.(в ред. Постановления Правительства РФ от 21.04.2016 N 335)
    \end{refsection}%
    \begin{refsection}[bl-author, bl-registered]
        % Это refsection=2.
        % Процитированные здесь работы:
        %  * попадают в авторскую библиографию, при usefootcite==0 и стиле `\insertbiblioauthorimportant`.
        %  * ни на что не влияют в противном случае
        % Примеры для стиля \insertbiblioauthorimportant. Впишите сюда ключи
        % своих работ из biblio/author.bib (пример --- Panov2025):
        \nocite{Panov2025}%
    \end{refsection}%
        %
        % Всё, что вне этих двух refsection, это refsection=0,
        %  * для диссертации - это нормальные ссылки, попадающие в обычную библиографию
        %  * для автореферата:
        %     * при usefootcite==0, ссылка корректно сработает только для источника из `external.bib`. Для своих работ --- напечатает "[0]" (и даже Warning не вылезет).
        %     * при usefootcite==1, ссылка сработает нормально. В авторской библиографии будут только процитированные в refsection=0 работы.
}

При использовании пакета \verb!biblatex! будут подсчитаны все работы, добавленные
в файл \verb!biblio/author.bib!. Для правильного подсчёта работ в~различных
системах цитирования требуется использовать поля:
\begin{itemize}
        \item \texttt{authorvak} если публикация индексирована ВАК,
        \item \texttt{authorscopus} если публикация индексирована Scopus,
        \item \texttt{authorwos} если публикация индексирована Web of Science,
        \item \texttt{authorconf} для докладов конференций,
        \item \texttt{authorpatent} для патентов,
        \item \texttt{authorprogram} для зарегистрированных программ для ЭВМ,
        \item \texttt{authorother} для других публикаций.
\end{itemize}
Для подсчёта используются счётчики:
\begin{itemize}
        \item \texttt{citeauthorvak} для работ, индексируемых ВАК,
        \item \texttt{citeauthorscopus} для работ, индексируемых Scopus,
        \item \texttt{citeauthorwos} для работ, индексируемых Web of Science,
        \item \texttt{citeauthorvakscopuswos} для работ, индексируемых одной из трёх баз,
        \item \texttt{citeauthorscopuswos} для работ, индексируемых Scopus или Web of~Science,
        \item \texttt{citeauthorconf} для докладов на конференциях,
        \item \texttt{citeauthorother} для остальных работ,
        \item \texttt{citeauthorpatent} для патентов,
        \item \texttt{citeauthorprogram} для зарегистрированных программ для ЭВМ,
        \item \texttt{citeauthor} для суммарного количества работ.
\end{itemize}
% Счётчик \texttt{citeexternal} используется для подсчёта процитированных публикаций;
% \texttt{citeregistered} "--- для подсчёта суммарного количества патентов и программ для ЭВМ.

Для добавления в список публикаций автора работ, которые не были процитированы в
автореферате, требуется их~перечислить с использованием команды \verb!\nocite! в
\verb!Synopsis/content.tex!.
%%%%%%%%%%%%%%%%%%%%%%%%%%%%%%%%%%%%%%%%%%%%%%%%%%%%%%%%%%%%%%%%%%%%%%%%%%%%%%%%%%
\fi
